% =====================================================================
%  Halloween-bursdagsinvitasjon for Sander
%  4 sider i A5 staaende (148,5 x 210 mm) = ett A4-ark liggende, brettet
%  loddrett pa midten. Leserekkefolge: 1 forside, 2-3 innside, 4 bakside.
%  Bygg med:  pdflatex invitasjon.tex   (kjor to ganger)
%  For utskrift: se invitasjon-utskrift.tex
% =====================================================================
\documentclass[11pt]{article}

\usepackage[T1]{fontenc}
\usepackage[utf8]{inputenc}
\usepackage[norsk]{babel}
\usepackage[sfdefault]{cabin}
\usepackage{comfortaa}
\usepackage{microtype}
\usepackage[paperwidth=148.5mm,paperheight=210mm,margin=0mm]{geometry}
\usepackage{tikz}
\usetikzlibrary{calc}

\renewcommand{\familydefault}{\sfdefault}
\pagestyle{empty}

% ---------- Opplysninger som er lette aa endre ------------------------
\newcommand{\bursdagsbarn}{Sander}
\newcommand{\adresse}{Solbakken 24}
\newcommand{\dagdato}{lørdag 31. oktober}
\newcommand{\starttid}{kl. 16.00}
\newcommand{\svarfrist}{så snart som mulig}
\newcommand{\kontakt}{Sigurd Rød Brekk, mobil 932 80 515}

% ---------- Farger ----------------------------------------------------
\definecolor{natt}{HTML}{2A1B3D}
\definecolor{nattlys}{HTML}{3E2A57}
\definecolor{krem}{HTML}{FFF4E2}
\definecolor{gresskar}{HTML}{F2892E}
\definecolor{gresskarm}{HTML}{DC6F1B}
\definecolor{gresskarl}{HTML}{FFA850}
\definecolor{stilk}{HTML}{5C8C3A}
\definecolor{maanefarge}{HTML}{FFE9A8}
\definecolor{skygge}{HTML}{1A1026}
\definecolor{lilla}{HTML}{7A5AA8}
\definecolor{rosa}{HTML}{FFB3C1}
\definecolor{bein}{HTML}{F7F1E3}
\definecolor{beinskygge}{HTML}{DDD2BA}

% ---------- Sidemal ---------------------------------------------------
% Hver side tegnes som ett fullt tikz-lerret med origo nede til venstre.
\newcommand{\side}[1]{%
  \begin{tikzpicture}[remember picture,overlay,x=1cm,y=1cm]
    \begin{scope}[shift={(current page.south west)}]
      #1
    \end{scope}
  \end{tikzpicture}\null\newpage}

\newcommand{\sidebredde}{14.85}
\newcommand{\sidehoyde}{21}
% Midten vannrett - brukes til aa sentrere alt paa sidene.
\newcommand{\midt}{7.425}

% ---------- Sikker sone ----------------------------------------------
%  De fleste kontorskrivere kan ikke printe helt ut i papirkanten -
%  typisk 4-8 mm forsvinner. Bakgrunnene blør derfor med vilje helt ut
%  (der gjør tapet ingenting, fargen er lik uansett), men ALT som betyr
%  noe - rammer, tekst, figurer - holdes minst \sikkersone fra kanten.
%  Da blir det hvit ramme rundt paa en vanlig skriver, men ingenting
%  blir beskaaret bort.
%
%  Flytter du paa noe: hold deg innenfor
%     x: 1,0 .. 13,85     y: 1,0 .. 20,0
\newcommand{\sikkersone}{1.0}

% Slaa av orddeling i loepende tekst (penere i korte avsnitt).
% NB: maa staa ETTER \selectfont, ellers gjelder den feil font.
\newcommand{\nohyph}{\hyphenchar\font=-1\relax}
\newcommand{\bakgrunn}[1]{\fill[#1] (0,0) rectangle (\sidebredde,\sidehoyde);}

% ---------- Tegninger -------------------------------------------------
% Gresskar: \gresskar[skala]{x}{y}
\newcommand{\gresskar}[3][1]{%
\begin{scope}[shift={(#2,#3)},scale=#1]
  \draw[stilk,line width=1.1mm,line cap=round]
        (0,0.9) .. controls (0.08,1.25) and (0.38,1.22) .. (0.30,1.55);
  \fill[gresskarm] (0,0) ellipse (1.25 and 1.00);
  \fill[gresskar]  (-0.52,0) ellipse (0.66 and 0.96);
  \fill[gresskar]  (0.52,0) ellipse (0.66 and 0.96);
  \fill[gresskarl] (0,0) ellipse (0.74 and 1.00);
  % oyne
  \fill[natt] (-0.40,0.26) -- (-0.10,0.26) -- (-0.25,0.60) -- cycle;
  \fill[natt] (0.40,0.26) -- (0.10,0.26) -- (0.25,0.60) -- cycle;
  % smil
  \fill[natt] (-0.52,0.02) .. controls (-0.26,-0.54) and (0.26,-0.54) .. (0.52,0.02)
              .. controls (0.26,-0.26) and (-0.26,-0.26) .. (-0.52,0.02) -- cycle;
  \fill[gresskarl] (-0.23,-0.19) -- (-0.07,-0.19) -- (-0.15,-0.40) -- cycle;
  \fill[gresskarl] (0.23,-0.19) -- (0.07,-0.19) -- (0.15,-0.40) -- cycle;
\end{scope}}

% Spokelse: \spokelse[skala]{x}{y}, kantfarge styres av \spokelsekant
\newcommand{\spokelsekant}{krem}
\newcommand{\spokelse}[3][1]{%
\begin{scope}[shift={(#2,#3)},scale=#1]
  \filldraw[fill=krem,draw=\spokelsekant,line width=0.35mm,line join=round]
    (-0.58,-0.45) -- (-0.58,0.12)
    .. controls (-0.58,0.98) and (0.58,0.98) .. (0.58,0.12) -- (0.58,-0.45)
    .. controls (0.53,-0.80) and (0.35,-0.80) .. (0.29,-0.45)
    .. controls (0.24,-0.80) and (0.06,-0.80) .. (0.00,-0.45)
    .. controls (-0.06,-0.80) and (-0.24,-0.80) .. (-0.29,-0.45)
    .. controls (-0.35,-0.80) and (-0.53,-0.80) .. (-0.58,-0.45) -- cycle;
  \fill[rosa] (-0.36,0.16) circle (0.08);
  \fill[rosa] (0.36,0.16) circle (0.08);
  \fill[natt] (-0.20,0.38) ellipse (0.08 and 0.11);
  \fill[natt] (0.20,0.38) ellipse (0.08 and 0.11);
  \fill[natt] (0,0.16) ellipse (0.09 and 0.07);
\end{scope}}

% Lite, snilt skjelett med skjerf og gresskarlykt: \skjelett[skala]{x}{y}
% (y = bakkenivaa, altsaa under foettene)
\newcommand{\skjelett}[3][1]{%
\begin{scope}[shift={(#2,#3)},scale=#1]
  % ---- skjaer fra lykta (bakerst)
  \shade[inner color=gresskarl,outer color=natt,opacity=0.55]
        (-0.72,0.28) circle (0.50);
  % ---- bein og foetter
  \draw[bein,line width=1.5mm,line cap=round] (-0.17,0.08) -- (-0.17,0.60);
  \draw[bein,line width=1.5mm,line cap=round] (0.17,0.08) -- (0.17,0.60);
  \fill[skygge,opacity=0.30] (0,0.03) ellipse (0.72 and 0.11);
  \fill[bein] (-0.23,0.07) ellipse (0.14 and 0.07);
  \fill[bein] (0.23,0.07) ellipse (0.14 and 0.07);
  % ---- hofte og brystkasse
  \fill[bein,rounded corners=0.07] (-0.25,0.56) rectangle (0.25,0.78);
  % brystkasse: smalere nede, med ribbein
  \fill[bein] (-0.34,1.38) .. controls (-0.42,1.10) and (-0.30,0.90) .. (-0.21,0.80)
              -- (0.21,0.80) .. controls (0.30,0.90) and (0.42,1.10) .. (0.34,1.38)
              -- cycle;
  \draw[beinskygge,line width=0.9mm,line cap=round] (0,1.34) -- (0,0.86);
  \foreach \yy/\dx in {1.24/0.29,1.08/0.31,0.92/0.26}{%
    \draw[beinskygge,line width=0.5mm,line cap=round]
      (-\dx,\yy+0.04) .. controls (0,\yy-0.07) .. (\dx,\yy+0.04);}
  % ---- armer: en ned med lykt, en vinker
  \draw[bein,line width=1.4mm,line cap=round,line join=round]
       (-0.35,1.28) -- (-0.62,1.04) -- (-0.72,0.76);
  \draw[bein,line width=1.4mm,line cap=round,line join=round]
       (0.35,1.30) -- (0.68,1.50) -- (0.86,1.80);
  \fill[bein] (-0.72,0.72) circle (0.10);
  \fill[bein] (0.88,1.84) circle (0.10);
  % ---- hodeskalle
  \draw[bein,line width=1.2mm,line cap=round] (0,1.34) -- (0,1.48);
  \fill[bein] (0,1.92) ellipse (0.44 and 0.46);
  \fill[bein] (0,1.62) ellipse (0.24 and 0.14);
  \fill[skygge] (-0.16,1.98) ellipse (0.115 and 0.135);
  \fill[skygge] (0.16,1.98) ellipse (0.115 and 0.135);
  \fill[bein] (-0.12,2.02) circle (0.035);
  \fill[bein] (0.12,2.02) circle (0.035);
  \fill[rosa,opacity=0.55] (-0.33,1.80) circle (0.09);
  \fill[rosa,opacity=0.55] (0.33,1.80) circle (0.09);
  \fill[skygge] (-0.05,1.80) -- (0.05,1.80) -- (0,1.71) -- cycle;
  \draw[skygge,line width=0.4mm,line cap=round]
       (-0.19,1.66) .. controls (-0.07,1.57) and (0.07,1.57) .. (0.19,1.66);
  \foreach \xx in {-0.09,0,0.09}{%
    \draw[skygge,line width=0.3mm,line cap=round] (\xx,1.60) -- (\xx,1.655);}
  % ---- skjerf
  \fill[gresskar,rounded corners=0.09] (-0.29,1.34) rectangle (0.29,1.51);
  \fill[gresskar] (0.03,1.51) -- (0.20,1.51) -- (0.24,1.14)
      .. controls (0.15,1.07) and (0.07,1.09) .. (0.05,1.16) -- cycle;
  \draw[krem,line width=0.3mm] (0.05,1.40) -- (0.22,1.40);
  \draw[krem,line width=0.3mm] (0.06,1.26) -- (0.23,1.26);
  % ---- gresskarlykt i haanden
  \gresskar[0.30]{-0.72}{0.26}
\end{scope}}

% Flaggermus: \flaggermus[skala]{x}{y}, farge styres av \flagfarge
\newcommand{\flagfarge}{skygge}
\newcommand{\flaggermus}[3][1]{%
\begin{scope}[shift={(#2,#3)},scale=#1]
  \fill[\flagfarge] (0,0) ellipse (0.15 and 0.19);
  \fill[\flagfarge] (-0.09,0.15) -- (-0.15,0.33) -- (-0.01,0.22) -- cycle;
  \fill[\flagfarge] (0.09,0.15) -- (0.15,0.33) -- (0.01,0.22) -- cycle;
  \fill[\flagfarge] (-0.11,0.09)
    .. controls (-0.42,0.30) and (-0.70,0.24) .. (-0.95,0.04)
    .. controls (-0.70,0.11) and (-0.58,-0.06) .. (-0.44,-0.13)
    .. controls (-0.30,-0.17) and (-0.19,-0.10) .. (-0.11,-0.05) -- cycle;
  \fill[\flagfarge] (0.11,0.09)
    .. controls (0.42,0.30) and (0.70,0.24) .. (0.95,0.04)
    .. controls (0.70,0.11) and (0.58,-0.06) .. (0.44,-0.13)
    .. controls (0.30,-0.17) and (0.19,-0.10) .. (0.11,-0.05) -- cycle;
  \fill[krem] (-0.06,0.03) circle (0.025);
  \fill[krem] (0.06,0.03) circle (0.025);
\end{scope}}

% Edderkopp i traad: \edderkopp[skala]{x}{y}
\newcommand{\edderkopp}[3][1]{%
\begin{scope}[shift={(#2,#3)},scale=#1]
  \draw[krem,opacity=0.55,line width=0.22mm] (0,0.25) -- (0,3.4);
  \foreach \s in {-1,1}{
    \draw[skygge,line width=0.5mm,line cap=round]
       (\s*0.16,0.08) .. controls (\s*0.45,0.26) .. (\s*0.58,-0.06);
    \draw[skygge,line width=0.5mm,line cap=round]
       (\s*0.18,0.02) .. controls (\s*0.50,0.10) .. (\s*0.62,-0.22);
    \draw[skygge,line width=0.5mm,line cap=round]
       (\s*0.17,-0.05) .. controls (\s*0.46,-0.06) .. (\s*0.56,-0.36);
    \draw[skygge,line width=0.5mm,line cap=round]
       (\s*0.14,-0.11) .. controls (\s*0.38,-0.20) .. (\s*0.44,-0.44);}
  \fill[skygge] (0,-0.05) circle (0.26);
  \fill[skygge] (0,0.20) circle (0.15);
  \fill[krem] (-0.06,0.24) circle (0.045);
  \fill[krem] (0.06,0.24) circle (0.045);
  \fill[skygge] (-0.06,0.24) circle (0.02);
  \fill[skygge] (0.06,0.24) circle (0.02);
\end{scope}}

% Stjerne: \stjerne[skala]{x}{y}
\newcommand{\stjerne}[3][1]{%
\begin{scope}[shift={(#2,#3)},scale=#1]
  \fill[krem,opacity=0.85] (0,0.13) .. controls (0.03,0.03) .. (0.13,0)
    .. controls (0.03,-0.03) .. (0,-0.13) .. controls (-0.03,-0.03) .. (-0.13,0)
    .. controls (-0.03,0.03) .. (0,0.13) -- cycle;
\end{scope}}

% Maane: \maane[skala]{x}{y}
\newcommand{\maane}[3][1]{%
\begin{scope}[shift={(#2,#3)},scale=#1]
  \fill[maanefarge,opacity=0.18] (0,0) circle (1.45);
  \fill[maanefarge] (0,0) circle (1.1);
  \fill[maanefarge!80!gresskar,opacity=0.5] (-0.35,0.30) circle (0.22);
  \fill[maanefarge!80!gresskar,opacity=0.5] (0.30,-0.10) circle (0.15);
  \fill[maanefarge!80!gresskar,opacity=0.5] (0.05,0.55) circle (0.10);
\end{scope}}

% Ramme til innsidene
\newcommand{\ramme}{%
  % Ligger innenfor \sikkersone, ellers kutter skriveren av rammen.
  \draw[gresskar,line width=0.9mm,rounded corners=6mm]
       (1.10,1.10) rectangle (13.75,19.90);
  \draw[gresskar,line width=0.25mm,rounded corners=5mm,opacity=0.6]
       (1.40,1.40) rectangle (13.45,19.60);}

\begin{document}

% =============================== SIDE 1: FORSIDE =====================
\side{%
  \bakgrunn{natt}
  \fill[nattlys,opacity=0.75] (\midt,0) ellipse (11.5 and 5.2);
  \maane[1]{12.3}{18.2}
  \foreach \p in {(1.7,19.2),(3.2,17.8),(5.0,19.4),(7.0,18.3),(9.2,19.5),
                  (1.9,15.6),(13.4,15.0),(1.4,12.9),(13.5,19.0),(4.2,15.2),
                  (13.4,12.2),(2.6,13.4)}
      {\stjerne[1]\p{}}
  \flaggermus[1.4]{9.1}{17.6}
  \flaggermus[0.95]{11.6}{15.7}
  \flaggermus[0.75]{3.1}{14.0}
  \edderkopp[1.0]{5.9}{15.9}

  \node[anchor=center,text=krem] at (\midt,13.5)
    {\fontsize{12}{15}\selectfont\textls[300]{DU ER INVITERT TIL}};
  \node[anchor=center,text=gresskarl] at (\midt,11.7)
    {\comfortaa\bfseries\fontsize{29}{34}\selectfont Halloween-bursdag};
  \node[anchor=center,text=krem] at (\midt,10.0)
    {\comfortaa\fontsize{22}{26}\selectfont hos \bursdagsbarn};

  \gresskar[1.9]{3.6}{6.6}
  \skjelett[1.6]{11.6}{4.70}

  % Stripa blør ut i kanten, men teksten i den holdes over \sikkersone.
  \fill[gresskar] (0,0) rectangle (\sidebredde,3.4);
  \node[anchor=center,text=natt] at (\midt,2.60)
    {\comfortaa\bfseries\fontsize{14}{17}\selectfont
     \dagdato\ \ $\cdot$\ \ \starttid};
  \node[anchor=center,text=natt] at (\midt,1.60)
    {\comfortaa\bfseries\fontsize{14}{17}\selectfont \adresse};
}

% =============================== SIDE 2: PROGRAM =====================
\side{%
  \bakgrunn{krem}
  \ramme
  \renewcommand{\flagfarge}{lilla}
  \renewcommand{\spokelsekant}{lilla}
  \flaggermus[0.7]{5.3}{19.0}
  \flaggermus[0.55]{12.7}{19.0}
  \flaggermus[0.5]{12.4}{2.7}

  \gresskar[0.7]{2.9}{18.1}
  \node[anchor=center,text=gresskarm] at (7.8,18.1)
    {\comfortaa\bfseries\fontsize{28}{32}\selectfont Program};
  \spokelse[0.75]{12.2}{18.0}

  \draw[gresskar,line width=0.4mm,dash pattern=on 0.6mm off 1.4mm,line cap=round]
       (2.60,16.0) -- (2.60,7.8);

  % Staaende format: klokkeslettet staar OVER teksten, ikke i egen
  % venstrespalte - spalta blir for smal naar sida er 14,85 cm bred.
  \foreach \y/\tid/\tekst in {%
    16.4/16.00/{Velkommen! Vi er ute, og bålet er tent.},
    13.4/{16.00--16.45}/{Utelek, og pølser som grilles på bålet.},
    10.4/16.45/{Vi går inn og varmer oss, og så spiser vi kake.},
     7.4/17.30/{Å gå knask-eller-knep i Solbakken for de som vil.}}{%
    \fill[gresskar] (2.60,\y) circle (0.16);
    \node[anchor=west,text=gresskarm] at (3.30,\y)
      {\comfortaa\bfseries\fontsize{16}{19}\selectfont kl. \tid};
    \node[anchor=north west,text=natt,align=left,text width=10.0cm] at (3.25,\y-0.5)
      {\fontsize{13.5}{17}\selectfont\nohyph \tekst};}

  \node[anchor=center,text=lilla,align=center,text width=11cm] at (\midt,4.6)
    {\comfortaa\fontsize{12}{16}\selectfont\nohyph
     Kle deg etter været, vi er mest ute!};
}

% =============================== SIDE 3: GODT AA VITE ================
\side{%
  \bakgrunn{krem}
  \ramme
  \renewcommand{\flagfarge}{lilla}
  \renewcommand{\spokelsekant}{lilla}
  \flaggermus[0.5]{12.6}{19.0}

  \node[anchor=center,text=gresskarm] at (\midt,18.5)
    {\comfortaa\bfseries\fontsize{28}{32}\selectfont Praktisk info};

  \foreach \y/\tekst in {%
    16.8/{Vi er mye ute, så kle deg etter været. Vi har tak over hodet ute hvis det regner.},
    14.8/{Utkledning er gøy, men helt frivillig. Kom akkurat som du vil.},
    12.8/{Vi serverer pølser som grilles på bålet, og kake inne.},
    10.8/{Naboene i Solbakken pynter mye, så runden kan bli litt skummel.},
     8.8/{Det er gratis parkering langs gaten utenfor.},
     6.8/{Er det noe vi bør vite om mat, allergier eller annet, så gi oss en lyd.}}{%
    \gresskar[0.30]{2.35}{\y}
    \node[anchor=north west,text=natt,align=left,text width=9.7cm] at (3.0,\y+0.45)
      {\fontsize{13}{17}\selectfont\nohyph \tekst};}

  \fill[gresskar,rounded corners=3mm] (1.95,3.1) rectangle (12.90,5.4);
  \node[anchor=center,text=natt,align=center,text width=10.2cm] at (7.425,4.25)
    {\comfortaa\fontsize{12}{17}\selectfont\nohyph
     Gi gjerne beskjed \svarfrist\\ til \kontakt};
}

% =============================== SIDE 4: BAKSIDE =====================
\side{%
  \bakgrunn{natt}
  \fill[nattlys,opacity=0.75] (\midt,0) ellipse (11 and 4.4);
  \foreach \p in {(1.8,19.0),(4.0,19.6),(6.6,18.4),(10.0,19.4),(12.6,18.7),
                  (13.4,15.6),(1.4,16.2),(13.3,11.6),(1.6,12.2),
                  (1.3,9.4),(13.4,9.0),(2.4,10.4),(12.4,13.4)}
      {\stjerne[0.9]\p{}}
  \flaggermus[1.05]{11.4}{17.2}
  \flaggermus[0.7]{3.0}{16.9}

  \node[anchor=center,text=krem,align=center,text width=12.4cm] at (\midt,13.6)
    {\comfortaa\fontsize{23}{29}\selectfont Vi gleder oss til\\ å feire med deg!};
  \node[anchor=center,text=gresskarl] at (\midt,10.7)
    {\comfortaa\bfseries\fontsize{20}{24}\selectfont Hilsen \bursdagsbarn};

  \gresskar[1.2]{3.3}{7.2}
  \spokelse[1.25]{11.6}{7.4}

  \node[anchor=center,text=krem] at (\midt,4.5)
    {\fontsize{13}{16}\selectfont \adresse};
  \node[anchor=center,text=krem] at (\midt,3.6)
    {\fontsize{13}{16}\selectfont \dagdato};
  \node[anchor=center,text=krem] at (\midt,2.7)
    {\fontsize{13}{16}\selectfont \starttid};
}

\end{document}
